% =====================================================================
% Appendix D -- Reproducibility
% =====================================================================
\documentclass[../preamble.tex]{subfiles}
\begin{document}

\chapter{Reproducibility}
\label{app:reproducibility}

\section{Hardware and software}
\label{app:repro-hw-sw}

\begin{table}[htbp]
  \centering
  \small
  \caption{Reproducibility summary.}
  \label{tab:repro-summary}
  \begin{tabular}{ll}
    \toprule
    \textbf{Item} & \textbf{Value} \\
    \midrule
    GPU            & NVIDIA RTX 4070 SUPER, 12\,GB \\
    OS             & Linux (Ubuntu 22.04 LTS) \\
    Python         & 3.10+ \\
    PyTorch        & 2.12.0+cu130 \\
    timm           & 1.0.27 \\
    scikit-learn   & 1.9.0 \\
    Optuna         & 4.9.0 \\
    scipy          & 1.17.1 \\
    numpy          & 2.4.4 \\
    OpenCV         & 4.13.0.92 \\
    matplotlib     & 3.10.9 \\
    Wall-clock (foundation + fine-tune) & approx.\ 5--6 hours \\
    \bottomrule
  \end{tabular}
\end{table}

\section{Seed policy}
\label{app:repro-seed}

The trainer sets Python, NumPy, and PyTorch seeds at the start of every
run. CUDA seeding is best-effort: deterministic algorithms are enabled
where available, and a deterministic fallback is documented for non-deterministic
operators.

\section{Repository entry points}
\label{app:repro-repo}

\begin{itemize}
  \item \texttt{scripts/generate\_analysis\_figures.py} -- reproduces the
        foundation-stage analysis figures under
        \texttt{docs/analysis/figures/}.
  \item \texttt{scripts/thesis\_figures.py} -- reproduces the thesis-specific
        figures and tables under \texttt{docs/thesis/}.
  \item \texttt{scripts/generate\_shortcut\_figures.py} -- reproduces the
        padding and border-attention figures.
  \item \texttt{scripts/generate\_shortcut\_figures\_nopad.py} -- the
        no-pad variant.
  \item \texttt{scripts/make\_paper\_figures.py} -- reproduces the paper
        figures under \texttt{docs/latex/figures/}.
\end{itemize}

The full fine-tune, HPO, and ensemble procedures are launched through the
internal training and evaluation entry points, whose YAML configurations are
checked into the repository. The thesis does not re-execute those procedures;
it consumes the artefacts they produced.

\section{Artefact inventory}
\label{app:repro-artifacts}

\subsection*{Headline metrics}
\path{results/finetune_zenodo/ensemble/top3_pre_hpo/external_validation/pcosgen.json}
\path{results/finetune_zenodo/ensemble/top3_pre_hpo/calibration/calibration_results.json}

\subsection*{Per-architecture fine-tune metrics}
\path{results/finetune_zenodo/checkpoints/srad_nopad/densenet121/final_metrics.json}
\path{results/finetune_zenodo/checkpoints/srad_nopad/convnext_tiny/final_metrics.json}
\path{results/finetune_zenodo/checkpoints/gauss_nopad/vit_base/final_metrics.json}

\subsection*{HPO best parameters}
\path{results/finetune_hpo/srad_nopad_densenet121/densenet121/best_params.yaml}
\path{results/finetune_hpo/srad_nopad_convnext_tiny/convnext_tiny/best_params.yaml}
\path{results/finetune_hpo/gauss_nopad_vit_base/vit_base/best_params.yaml}

\subsection*{MC-Dropout predictions}
\path{results/finetune_zenodo/checkpoints/srad_nopad/densenet121/uncertainty/mc_dropout_predictions.csv}
\path{results/finetune_zenodo/checkpoints/srad_nopad/convnext_tiny/uncertainty/mc_dropout_predictions.csv}
\path{results/finetune_zenodo/checkpoints/gauss_nopad/vit_base/uncertainty/mc_dropout_predictions.csv}

\subsection*{External validation predictions (per-checkpoint)}
\path{results/finetune_zenodo/checkpoints/srad_nopad/densenet121/external_validation/pcosgen.csv}
\path{results/finetune_zenodo/checkpoints/srad_nopad/convnext_tiny/external_validation/pcosgen.csv}
\path{results/finetune_zenodo/checkpoints/gauss_nopad/vit_base/external_validation/pcosgen.csv}

\subsection*{Sweep matrices}
\path{results/ablation/sweep_matrix.csv}
\path{results/ablation_nopad/sweep_matrix.csv}

\section{Verifying headline numbers}
\label{app:repro-verify}

A reader can verify any headline number in this thesis by loading the
corresponding JSON or CSV and checking it against the reported value. The
script below reproduces every headline metric in
\cref{tab:headline-numbers}:

\begin{verbatim}
import json
ens = json.load(open(
    "results/finetune_zenodo/ensemble/top3_pre_hpo/"
    "external_validation/pcosgen.json"))
cal = json.load(open(
    "results/finetune_zenodo/ensemble/top3_pre_hpo/"
    "calibration/calibration_results.json"))
assert abs(ens["test_auc_roc"] - 0.9487) < 1e-4
assert abs(ens["test_f1"] - 0.8665) < 1e-4
assert abs(ens["test_mcc"] - 0.7884) < 1e-4
assert abs(cal["ece_after_pass2"] - 0.0810) < 1e-4
print("all headline metrics match")
\end{verbatim}

\section{Wall-clock note}
\label{app:repro-wallclock}

The wall-clock estimate is the sum of the time-stamped training logs in the
sweep matrices, plus the HPO trials, the post-HPO retrain, the two-pass
calibration, the MC-Dropout uncertainty evaluation, and the Grad-CAM panel
generation. The reported 5--6 hours is conservative. A user with the same
hardware could expect a similar total; a different GPU, batch size, or
PyTorch version would change it.

\end{document}